\documentclass[aspectratio=169]{beamer}
\input{header}
\coursecode{JKSSB}

\title{Access Queries, Forms \& Reports}
\subtitle{Lesson 30 --- Objects, Queries and Reporting}
\author{JKSSB Computer Awareness}
\date{\today}

\begin{document}
\frame{\titlepage}

\begin{frame}{The Four Main Access Objects}
  \begin{center}
    \begin{tabular}{p{0.16\textwidth}p{0.72\textwidth}}
      \toprule
      \textbf{Object} & \textbf{Job} \\
      \midrule
      Table & Stores the actual data in rows and columns \\
      Query & Asks a question and retrieves selected data \\
      Form & Provides a screen to enter and view data \\
      Report & Presents data in a printable layout \\
      \bottomrule
    \end{tabular}
  \end{center}
  \vspace{0.25cm}
  \muted{Data lives in tables; queries, forms and reports work on that data.}
\end{frame}

\begin{frame}{The Object Sequence}
  \begin{center}
    \key{Table} $\rightarrow$ \key{Query} $\rightarrow$ \key{Form}
    $\rightarrow$ \key{Report}
  \end{center}
  \vspace{0.2cm}
  \begin{itemize}
    \item \key{Table}: the store where the data actually lives.
    \item \key{Query}: asks a question and pulls the needed rows together.
    \item \key{Form}: a screen for entering and viewing data.
    \item \key{Report}: a formatted, printable presentation of the data.
  \end{itemize}
  \begin{block}{The one idea}
    Data is stored once, in a table; the query selects it, the form enters
    it, and the report prints it.
  \end{block}
\end{frame}

\begin{frame}{Tables Revisited: The Source of Data}
  \begin{itemize}
    \item A \key{table} stores data in \key{fields} (columns) and
      \key{records} (rows).
    \item Every other object draws its data from a \key{table} or a
      \key{query}.
    \item A form or a report built on a table shows that table's data; built
      on a query, it shows only the \key{selected} rows.
    \item The table is therefore the \key{foundation} of the database.
  \end{itemize}
  \begin{block}{Remember}
    No data is stored in a query, form or report --- they only
    \muted{present} what a table holds.
  \end{block}
\end{frame}

\begin{frame}{What Is a Query?}
  \begin{itemize}
    \item A \key{query} is a question you ask of the database.
    \item It \key{retrieves} data from one or more tables, usually by a
      condition.
    \item A query does not normally store new data; it shows a result
      \key{set} drawn from the tables.
    \item Queries can also \key{modify} data --- these are called
      \key{action queries}.
  \end{itemize}
  \begin{block}{Definition}
    A query selects, filters, or acts on records according to the rules you
    set.
  \end{block}
\end{frame}

\begin{frame}{The Select Query}
  \begin{itemize}
    \item A \key{Select query} \key{retrieves} records that satisfy a
      \key{condition} (a criterion).
    \item It is the \key{most common} query type in Access.
    \item Example: show every student whose \texttt{Marks} are greater than
      60.
    \item The condition decides \key{which rows appear}; the fields decide
      \key{which columns appear}.
  \end{itemize}
  \begin{alertblock}{Exam point (PYQ-39)}
    A query that retrieves data on the basis of a condition is a
    \key{Select Query}.
  \end{alertblock}
\end{frame}

\begin{frame}{Building a Query: Design View and Criteria}
  \begin{itemize}
    \item Queries are built in \key{Query Design} view (or typed in
      \key{SQL} view).
    \item The design grid is called \key{QBE} --- \key{Query By Example}.
    \item Each column holds a \key{field}; the \key{Criteria} row holds the
      condition.
    \item Example criterion: \texttt{>60} on the \texttt{Marks} field shows
      only rows above 60.
  \end{itemize}
  \begin{exampleblock}{Worked example}
    For a Students table, put \texttt{Marks} in the field row and
    \texttt{>60} in the Criteria row; the query returns only students who
    scored more than 60.
  \end{exampleblock}
\end{frame}

\begin{frame}{Action Queries: Update, Append, Delete}
  \begin{center}
    \begin{tabular}{p{0.18\textwidth}p{0.68\textwidth}}
      \toprule
      \textbf{Query} & \textbf{What it does to the data} \\
      \midrule
      Update & Changes existing values in records \\
      Append & Adds records from one table to another \\
      Delete & Removes records that match a condition \\
      \bottomrule
    \end{tabular}
  \end{center}
  \vspace{0.2cm}
  \muted{Update, Append and Delete all change the stored data; that is why
  they are called action queries.}
\end{frame}

\begin{frame}{Select vs Action Queries (Near-Neighbour)}
  \begin{center}
    \begin{tabular}{p{0.20\textwidth}p{0.34\textwidth}p{0.32\textwidth}}
      \toprule
      \textbf{Point} & \textbf{Select query} & \textbf{Action query} \\
      \midrule
      What it does & Retrieves data on a condition & Modifies stored data \\
      Changes data? & No & Yes \\
      Examples & Select & Update, Append, Delete \\
      \bottomrule
    \end{tabular}
  \end{center}
  \vspace{0.25cm}
  \begin{alertblock}{Trap alert (TRAP-31)}
    A \key{Select} query only \key{retrieves}; to modify data use an
    \key{Update}, \key{Append}, or \key{Delete} query.
  \end{alertblock}
\end{frame}

\begin{frame}{Lookup Fields}
  \begin{itemize}
    \item A \key{lookup} field lets the user \key{pick a value from a list}
      instead of typing it.
    \item The list usually comes from a \key{table} or a \key{query}, such
      as a list of department names.
    \item A lookup is often used for a \key{foreign key}, so the entry
      matches an existing value.
    \item Benefit: it \key{reduces typing errors} and keeps values
      consistent.
  \end{itemize}
  \begin{block}{Keep it clear}
    A lookup field \emph{chooses} a value from a list; it does not create a
    new table or a new query.
  \end{block}
\end{frame}

\begin{frame}{Basic SQL Awareness}
  \begin{itemize}
    \item \key{SQL} stands for \key{Structured Query Language}.
    \item It is the standard language for working with relational databases;
      Access runs SQL behind the scenes.
    \item A typical retrieval statement is
      \texttt{SELECT fields FROM table WHERE condition;}.
    \item \texttt{SELECT} retrieves; \texttt{FROM} names the table;
      \texttt{WHERE} gives the condition.
  \end{itemize}
  \begin{exampleblock}{Example}
    \texttt{SELECT Name FROM Students WHERE Marks > 60;} returns the
    \texttt{Name} of every student with marks above 60.
  \end{exampleblock}
\end{frame}

\begin{frame}{Forms: Entering and Viewing Data}
  \begin{itemize}
    \item A \key{form} is an object that provides a \key{screen} --- a user
      interface --- for entering, viewing, and editing data.
    \item A form is usually \key{based on a table or a query}.
    \item It shows one record at a time (or a few), with clearly labelled
      fields.
    \item Forms make data entry \key{easier and safer} than typing directly
      into a table grid.
  \end{itemize}
  \begin{alertblock}{Exam point}
    A form is for \key{on-screen} work with data --- entering and viewing it.
  \end{alertblock}
\end{frame}

\begin{frame}{Reports: Printable Output}
  \begin{itemize}
    \item A \key{report} presents data in a \key{formatted, printable}
      layout.
    \item It can \key{group} records and show \key{totals and summaries}.
    \item A report is \key{read-only}: it is meant for \key{printing and
      sharing}, not for editing data.
    \item A report can be based on a table or a query.
  \end{itemize}
  \begin{block}{Definition}
    A report is the print-ready view of your data --- the output object.
  \end{block}
\end{frame}

\begin{frame}{Form vs Report (Near-Neighbour)}
  \begin{center}
    \begin{tabular}{p{0.18\textwidth}p{0.34\textwidth}p{0.34\textwidth}}
      \toprule
      \textbf{Point} & \textbf{Form} & \textbf{Report} \\
      \midrule
      Purpose & Enter and view data on screen & Print and present data \\
      Data change & Allows editing & Read-only \\
      Typical use & Data-entry screen & Printed summary \\
      \bottomrule
    \end{tabular}
  \end{center}
  \vspace{0.25cm}
  \muted{Form = screen for data entry; Report = printed output. Do not swap
  them.}
\end{frame}

\begin{frame}{Worked Example: A Students Database}
  \begin{itemize}
    \item \key{Table}: \texttt{Students} with \texttt{RollNo},
      \texttt{Name}, \texttt{Marks}.
    \item \key{Query}: select students with \texttt{Marks} greater than 60
      (a Select query).
    \item \key{Form}: a screen to enter a new student's details.
    \item \key{Report}: a printable list of all students who passed.
  \end{itemize}
  \begin{exampleblock}{The full sequence}
    Table (stores the data) $\rightarrow$ Query (selects it) $\rightarrow$
    Form (enters it) $\rightarrow$ Report (prints it).
  \end{exampleblock}
\end{frame}

\begin{frame}{Trap Alert: Objects and Query Types}
  \begin{itemize}
    \item A \key{Select} query \key{retrieves}; \key{Update}, \key{Append}
      and \key{Delete} queries \key{modify} (TRAP-31).
    \item Data is stored in a \key{table}, \key{not} in a query, form or
      report.
    \item A \key{form} is for \key{on-screen entry}; a \key{report} is a
      \key{printable} output.
    \item \key{SQL} = Structured Query Language.
    \item A \key{lookup} field picks a value \key{from a list}.
  \end{itemize}
\end{frame}

\begin{frame}{Lesson 30: Recall}
  \begin{itemize}
    \item The four main objects: \key{Table} (stores), \key{Query}
      (retrieves), \key{Form} (enters on screen), \key{Report} (prints).
    \item The sequence is \key{Table} $\rightarrow$ \key{Query}
      $\rightarrow$ \key{Form} $\rightarrow$ \key{Report}.
    \item A \key{Select query} retrieves data on a condition; \key{action
      queries} (Update, Append, Delete) modify data.
    \item Queries are built in \key{QBE} (Query By Example) design view with
      a \key{Criteria} row.
    \item A \key{lookup} field chooses a value from a list.
    \item \key{SQL} (Structured Query Language) uses
      \texttt{SELECT ... FROM ... WHERE}.
  \end{itemize}
\end{frame}
\end{document}
